\begin{trapbox}
	\begin{description}[style=nextline, leftmargin=2.8em, labelsep=0.5em, itemsep=0.35em]
		\item[\textbf{\textcolor{CTrap}{易错点一（忽略平行前提）}}] 误认为只要一个平面与两个平面相交，所得交线就一定平行，而忽略了“两个平面必须平行”这一关键条件。
		\item[\textbf{\textcolor{CTrap}{正确理解（平行前提必要）：}}] 只有两个平面平行时，第三个平面与它们的交线才平行；若两平面不平行，交线可能相交或异面。
		\item[\textbf{\textcolor{CTrap}{错因分析（前提记忆缺失）：}}] 只记住了结论的结果，没有记住定理的条件，导致乱用。
	\end{description}

	\medskip
	\begin{description}[style=nextline, leftmargin=2.8em, labelsep=0.5em, itemsep=0.35em]
		\item[\textbf{\textcolor{CTrap}{易错点二（误认异面直线）}}] 误以为两条交线是异面直线，忽略它们共面于第三平面，从而不能直接判断平行。
		\item[\textbf{\textcolor{CTrap}{正确理解（共面性质）：}}] 两条交线都位于第三平面内，因此一定共面；结合两平面平行可证它们不相交，从而平行。
		\item[\textbf{\textcolor{CTrap}{错因分析（共面意识不足）：}}] 空间想象力弱，没有意识到交线属于第三平面内的直线。
	\end{description}

	\medskip
	\begin{description}[style=nextline, leftmargin=2.8em, labelsep=0.5em, itemsep=0.35em]
		\item[\textbf{\textcolor{CTrap}{易错点三（交线识别错误）}}] 在复杂几何体中，找不准哪两条线段是同一个平面与两个平行平面的交线，从而无法应用定理。
		\item[\textbf{\textcolor{CTrap}{正确理解（识别对应交线）：}}] 应当先明确第三个平面是什么，然后分别找出它与两个平行面的交线，这些交线通常位于几何体的边界上。
		\item[\textbf{\textcolor{CTrap}{错因分析（图形辨识不足）：}}] 未能将实物图转化为抽象的交线模型，缺乏训练。
	\end{description}
\end{trapbox}
